\chapter{High Bandwidth Flash: A Technical Dossier}
\chaptermark{High Bandwidth Flash}
\label{app:hbf}

Section~\ref{sec:memoryunlock} rests part of its argument on a memory technology that does not yet exist as a shipping product, and Section~\ref{sec:appliance-cost} prices a machine built around it. A book that does that owes the reader a full account of what is actually known about the technology, what is claimed, what is modelled, and what is not published at all---in one place, graded, and dated. This appendix is that account, current as of \datafreeze. It is organized so that a reader can check the load-bearing assumptions of Chapter~\ref{ch:compute} against the primary record: the lineage of the idea, the announcement history, the specification as published, the power and endurance physics, the cost structure, the perceived advantages and their strongest rebuttals, the competitive landscape, the China angle, and the research literature. The short version is that HBF is real as a specification, absent as silicon, plausible as a weights tier, implausible as anything else, and priced---by its vendors---in adjectives rather than numbers.

\section{What it is, and why this book cares}

High Bandwidth Flash is NAND flash reorganized to behave like a capacity-first cousin of High Bandwidth Memory: sixteen (or eight) NAND dies stacked with through-silicon vias on a logic base die, in a package that matches HBM4's footprint and stack height, read in parallel across many independent channels so that a single 512\,GB stack delivers on the order of a terabyte and a half per second---HBM3E-class read bandwidth at eight to sixteen times HBM's capacity---and connected to the host processor over UCIe rather than PCIe~\cite{sandisk-hbf-facts,hbf-ocp,hbf-spec}. Removing the SSD controller and the PCIe link is what turns flash from a storage device into a memory tier: a PCIe~5.0 SSD delivers about 14\,GB/s and an NVMe path costs on the order of a hundred picojoules per bit read~\cite{ssd-energy}, whereas the HBF proposal is a hundred times the bandwidth at what its proponents model as a tenth of the energy per bit. It is non-volatile, so it draws no refresh power and holds a model across power cycles. It is slow to write and wears when written, so it is proposed for weights and read-mostly context, not for key--value cache or activations. That is the entire concept, and every argument about it is an argument about whether those five properties---capacity, bandwidth, energy, non-volatility, write asymmetry---land where the vendors say they will.

The book cares because a terabyte-class weights tier at flash cost is the specific mechanism by which Section~\ref{sec:memoryunlock} argues the frontier open-weight commons could move from the datacentre to the server closet, and because the appliance priced in Section~\ref{sec:appliance-cost} carries four such stacks. If HBF ships at its specification, that machine is buildable; if it ships at half its bandwidth or twice its energy, the machine is slower or hotter but still viable; if it does not ship, the argument falls back to LPDDR-only capacity tiers, which Appendix~\ref{app:lpddr} shows are adequate but roughly three times costlier per gigabyte. The dependence is therefore real but not total, and this appendix is careful to say which conclusions of the main text hang on which unverified numbers.

\section{Lineage: three earlier attempts to make flash behave like memory}

HBF is the fourth serious attempt to put non-volatile memory into the latency-and-bandwidth band between DRAM and storage, and its designers have said as much. Toshiba Memory (now Kioxia) introduced XL-FLASH, a single-level-cell storage-class memory with roughly five-microsecond read latency, at the Flash Memory Summit of August 2019, and it remains in production as the fast tier of Kioxia's enterprise SSDs and, since 2026, of a CXL memory module~\cite{hbf-xlflash} [V]. Intel's Optane---3D~XPoint, the most ambitious of the attempts and the only one to ship as a DIMM---was wound down in July 2022 with a \$559M inventory impairment, having failed to find a market large enough to fund a proprietary process at two vendors~\cite{hbf-optane} [V]. Samsung revived its Z-NAND at FMS 2025 as a GPU-direct storage tier, claiming fifteenfold the performance and 80\% lower power than conventional NAND~\cite{hbf-znand} [A]. The lesson the industry drew from Optane, and which every critic of HBF cites, is that a memory technology owned by one or two vendors and requiring a new controller, a new interface and a new software tier faces an adoption problem that its physics cannot solve; the lesson HBF's proponents draw is different---that Optane failed on \emph{cost per bit}, where it sat awkwardly between DRAM and NAND, whereas HBF is built from the cheapest bits in the industry and asks only for packaging that HBM has already made routine~\cite{hbf-optane-analogy} [A].

What is new in HBF is not the flash but the assembly. Three enabling pieces arrived together in 2024--25: SanDisk's CBA (``CMOS directly Bonded to Array'') architecture, which places the peripheral logic on a separate wafer bonded to the memory array and thereby frees the array die for the many independent, concurrently addressable sub-arrays that high aggregate read bandwidth requires; through-silicon vias through NAND dies, which HBM had made an industrial process for DRAM but which had not previously been productized for flash; and the UCIe die-to-die standard, which gave a package-internal interface that neither HBM's PHY nor PCIe provided~\cite{sandisk-hbf-blog,hbf-semieng} [V/A]. A fourth piece, the sixty-fold gap between what an HBM stack and a NAND stack cost per gigabyte, is what makes the assembly worth doing.

\section{The announcement record, February 2025 to August 2026}

Table~\ref{tab:hbf-timeline} sets out the public record. Three features of it deserve comment before the specification itself. First, the timeline has slipped in the way that first-of-a-kind memory timelines slip: in August 2025 the vendors said samples in the second half of 2026 and first inference devices in early 2027; by April 2026 SanDisk's chief executive was saying ``the NAND late this year, and \ldots more of a system with the controller early to mid next year''; by the August 2026 investor day the deck said ``first HBF inference product samples 2027,'' and the sell-side was writing 2028 for production~\cite{sandisk-hbf-blog,hbf-q3call,hbf-investorday,hbf-tapeout} [V/A]. Second, the coalition has widened but not to the parties that matter most: Google (through DeepMind, whose Xiaoyu Ma co-authored the Patterson paper that gives HBF its academic case) and Tenstorrent joined the OCP consortium in August 2026, and Meta was reported to have joined at the investor day; NVIDIA, AMD, Samsung, Micron and Kioxia have made no public commitment, and Kioxia---SanDisk's own fab partner---is pursuing a separate high-bandwidth-flash path with NVIDIA~\cite{hbf-ocp,hbf-investorday,hbf-kioxia-tsv,hbf-toms-spec} [V/A]. Third, the one hard milestone is recent and unambiguous: SanDisk stated on 13~August 2026 that the first HBF memory die has taped out~\cite{hbf-investorday,hbf-tapeout} [V]. Silicon exists; product does not.

\begin{table}[htbp]
  \centering
  \scriptsize
  \caption{High Bandwidth Flash: the announcement record. V = verified against the primary page; A = trade or analyst reporting; R = roadmap statement.}
  \label{tab:hbf-timeline}
  \begin{tabular}{p{16mm}p{57mm}p{62mm}p{7mm}}
    \toprule
    Date & Event & What was said & Grade \\
    \midrule
    Aug.\ 2019 & Toshiba/Kioxia XL-FLASH at FMS & SLC storage-class memory, $\sim$5\,\textmu s latency; the flash-as-memory precedent~\cite{hbf-xlflash} & V \\
    Jul.\ 2022 & Intel winds down Optane & \$559M impairment; the two-vendor cautionary tale~\cite{hbf-optane} & V \\
    11 Feb.\ 2025 & SanDisk Investor Day: first HBF disclosure & ``Match the bandwidth of HBM \ldots 8 to 16 times capacity at a similar cost point''; 512\,GB per 16-high stack; ``4\,TB of VRAM'' on a GPU~\cite{hbf-toms-2025,hbf-bf-2025} & V/A \\
    Jul.\ 2025 & SanDisk HBF Fact Sheet & Gen.\ 1: 512\,GB, 1.6\,TB/s; Gen.\ 2: up to 1\,TB, $>$2\,TB/s, 0.8$\times$ power; Gen.\ 3: 1.5\,TB, 3.2\,TB/s, 0.64$\times$; ``non-volatile, no refresh power needed''~\cite{sandisk-hbf-facts} & V \\
    24 Jul.\ 2025 & HBF Technical Advisory Board & David Patterson (chair) and Raja Koduri; Jim Keller added by Aug.\ 2026~\cite{hbf-tab,hbf-investorday} & V \\
    6 Aug.\ 2025 & FMS 2025: SanDisk--SK hynix MoU & ``establish specifications''; samples 2H CY2026, first inference devices early 2027; Best of Show~\cite{hbf-mou,sandisk-hbf-blog} & V \\
    20 Aug.\ 2025 & Kioxia 5\,TB / 64\,GB/s flash module & PCIe~6.0 daisy-chained ``beads,'' $<$40\,W, NEDO-funded; not HBF, the alternative path~\cite{hbf-kioxia-module} & V \\
    Oct.\ 2025 & OCP Global Summit & SK hynix ``AIN B'' product family ``leverages HBF technology''; joint ``HBF Night''~\cite{hbf-skh-ocp25} & V \\
    Nov.\ 2025 & Trade analysis & ``two or more years from commercial availability''; NVIDIA involvement ``essential''~\cite{hbf-bf-years} & A \\
    Dec.\ 2025 & YMTC at its developer conference & HBF fits Xtacking; ``full speed toward HBF next'' (DigiTimes)~\cite{hbf-ymtc} & A \\
    8 Jan.\ 2026 & Ma \& Patterson, arXiv:2601.05047 & HBF as ``10$\times$ memory capacity with HBM-like bandwidth''; ballpark stack 512\,GB, 1.6\,TB/s, $<$80\,W~\cite{hbf-ma-patterson} & V \\
    Jan.\ 2026 & SK hynix H$^3$ paper (IEEE CAL) & HBM+HBF hybrid: 2.69$\times$ perf/W, 18.8$\times$ batch at 10M-token context; assumes $\sim$100k write cycles, up to 4$\times$ HBM power~\cite{hbf-h3} & V \\
    25 Feb.\ 2026 & OCP consortium launched & ``dedicated workstream under OCP''; SanDisk and SK hynix founders~\cite{hbf-consortium} & V \\
    20 Mar.\ 2026 & Kioxia's own path & up to 32 TSV-stacked dies on a shared interposer beside HBM, ``with Nvidia''; XL-Flash the candidate cell~\cite{hbf-kioxia-tsv} & A \\
    30 Apr.\ 2026 & SanDisk Q3 FY26 call & ``the NAND late this year \ldots a system with the controller early to mid next year''; ``not a substitute for an enterprise SSD''~\cite{hbf-q3call} & A \\
    Jun.\ 2026 & Equipment and IP & Hanmi first HBF TC bonders 2H26; SanDisk patent bonding a processor onto a CBA NAND tile with HBM on a shared interposer~\cite{hbf-equipment} & A \\
    28 Jun.\ 2026 & Micron at HotInfra (ISCA) & roofline: on-package LPDDR beats HBF at every operating point for datacentre decode; HBF 8.5\,pJ/bit, 20\,\textmu s~\cite{hbf-micron-hotinfra} & V \\
    11 Jul.\ 2026 & FlashAccel, arXiv:2607.10186 (CAS) & 6 HBF stacks + HBM: 2.54$\times$ throughput per GPU, 1.93$\times$ energy efficiency at a 100\,ms SLO~\cite{hbf-flashaccel} & V \\
    3--5 Aug.\ 2026 & OCP specification v0.7.0; FMS 2026 & 8-Hi/16-Hi to 512\,GB; three grades $\sim$0.4/1.5/3.0\,TB/s; UCIe; Google and Tenstorrent join~\cite{hbf-ocp,hbf-spec,hbf-bf-spec} & V \\
    5 Aug.\ 2026 & Samsung and Kioxia at FMS 2026 & Samsung zNAND-O (4/8-die NAND stack for on-device NPU) and zHBM; Kioxia stays out of OCP HBF~\cite{hbf-samsung-fms26} & V/A \\
    13 Aug.\ 2026 & SanDisk Investor Day & ``First HBF memory die taped out''; ``first HBF inference product samples 2027''; ``4 HBF GPUs = 8 HBM GPUs''~\cite{hbf-investorday} & V \\
    14 Aug.\ 2026 & Sell-side and critics & Goldman, Cantor bullish; Counterpoint/Wedbush: unmodelled upside; critics: HBM4E baseline, BF16 assumption, no endurance figure~\cite{hbf-tapeout,hbf-counterpoint,hbf-xeno} & A \\
    \bottomrule
  \end{tabular}
\end{table}

\section{The specification, as published, and what it omits}

The document released on 3~August 2026 is the \emph{OCP HBF Architecture Specification, Version 0.7.0}, subtitled a high-level base-die specification---not a 1.0, and developed, by SanDisk's account, in about six months inside an OCP storage workstream chosen over JEDEC because its ``workstreams are highly goal-oriented'' and can ``iterate on the specification in real-time''~\cite{hbf-bf-spec,hbf-semieng} [V/A]. What it fixes, per the vendors' own summaries and the trade outlets that read it, is set out in Table~\ref{tab:hbf-spec}, and what it does not fix is the more important half of the table.

\begin{table}[htbp]
  \centering
  \small
  \caption{HBF as specified and as claimed. ``Published'' means stated by SanDisk, SK hynix or the OCP document; ``modelled'' means a third-party assumption; ``none'' means no figure of any kind has been published as of \datafreeze.}
  \label{tab:hbf-spec}
  \begin{tabular}{p{36mm}p{56mm}p{20mm}p{30mm}}
    \toprule
    Quantity & Value & Status & Source \\
    \midrule
    Capacity per stack & up to 512\,GB (16 $\times$ 256\,Gb dies); 8-Hi and 16-Hi & published & \cite{hbf-spec,sandisk-hbf-facts} \\
    Roadmap capacity & Gen.\ 2 ``up to 1\,TB,'' Gen.\ 3 1.5\,TB (fact sheet); 1.5$\times$ / 2$\times$ (SemiEngineering) & published, inconsistent & \cite{sandisk-hbf-facts,hbf-semieng} \\
    Read bandwidth grades & $\sim$0.4 / 1.5 / 3.0\,TB/s (0.384 / 1.536 / 3.072); Gen.\ 1 product 1.6\,TB/s & published & \cite{hbf-spec,hbf-bf-spec} \\
    Host interface & UCIe, reported up to 64\,GT/s over 64 lanes ($\sim$400--500\,GB/s per module) & published / reported & \cite{hbf-toms-spec} \\
    Base die & up to 16 independent full-duplex host channels, 1--4 AXI channels each (up to 64); each host channel owns a set of NAND dies & reported from spec & \cite{hbf-bf-spec} \\
    Package & ``closely matching the physical footprint, power profile, and stack height of HBM4''; no dimensions & published, qualitative & \cite{sandisk-hbf-facts} \\
    Interposer & not stated; concepts show HBM and HBF stacks side by side on one interposer & modelled & \cite{hbf-kioxia-tsv,hbf-h3} \\
    Cell type / node & not published; sources say SLC-like; BiCS generation undisclosed & none & \cite{hbf-analysis,hbf-bf-spec} \\
    Read latency & ``higher than HBM''; microsecond range; modelled at 4\,\textmu s (SLC) to 20\,\textmu s & none from vendor & \cite{hbf-flashaccel,hbf-micron-hotinfra} \\
    Read granularity & ``larger page size than HBM''; modelled $\sim$128$\times$ HBM4; 4\,KB pages & none from vendor & \cite{hbf-analysis,hbf-haven} \\
    Write bandwidth & supported (``read/write user guide''); no figure; NAND program $\sim$75--85\,MB/s per die & none & \cite{hbf-spec,hbf-isscc} \\
    Power & ``similar power profile'' to HBM4; Gen.\ 2 0.8$\times$, Gen.\ 3 0.64$\times$ Gen.\ 1; no watts, no pJ/bit & relative only & \cite{sandisk-hbf-facts} \\
    Modelled energy per bit & 6 (Ma--Patterson), 8 (FlashAccel), 8.5 (Micron) pJ/bit; HBM 3.0--4.2 & modelled & \cite{hbf-ma-patterson,hbf-flashaccel,hbf-micron-hotinfra,ssd-energy} \\
    Standby & ``non-volatile, no refresh power needed'' & published & \cite{sandisk-hbf-facts} \\
    Endurance & no P/E, TBW or read-disturb figure; H$^3$ assumes $\sim$100k cycles, FlashAccel 1M & none & \cite{hbf-h3,hbf-flashaccel} \\
    Retention at package temperature & no figure; 3D NAND retention loss accelerates superlinearly with temperature & none & \cite{hbf-heatwatch} \\
    ECC / reliability scheme & spec has a ``reliability'' section; contents not public & none & \cite{hbf-ocp} \\
    Stack price & ``similar cost'' to an HBM stack; ``one of the lowest cost per bit'' & adjectives & \cite{sandisk-hbf-facts,hbf-toms-2025} \\
    Samples / production & first die taped out Aug.\ 2026; product samples 2027; production 2028 (sell-side) & published / A & \cite{hbf-investorday,hbf-tapeout} \\
    \bottomrule
  \end{tabular}
\end{table}

Two things in that table matter more than the rest. The first is that every number a system designer would need to commit to a design---latency, granularity, write bandwidth, watts, endurance, retention, price---is either absent or someone else's assumption; the vendors have published capacity, bandwidth grades and an interface, and described everything else in comparatives. Contemporaneous critics made the same point more sharply: the OCP document carries no version discipline beyond ``0.7.0,'' no conformance procedure and no certification scheme, and its bandwidth and capacity figures are two vendors' development targets rather than requirements~\cite{hbf-critique} [A]. The second is that the roadmap itself is stated inconsistently between SanDisk's own documents---the July 2025 fact sheet says Gen.\ 2 reaches 1\,TB and Gen.\ 3 reaches 1.5\,TB, while the roadmap SanDisk gave Semiconductor Engineering in May 2026 says 1.5$\times$ and 2$\times$ Gen.\ 1 capacity, i.e.\ 768\,GB and 1\,TB~\cite{sandisk-hbf-facts,hbf-semieng}. Section~\ref{sec:memoryunlock} quotes the fact sheet; the reader should know that the more recent statement is the more conservative one, and that the difference is a full generation of capacity.

\section{Power, energy, and the physics of reading flash fast}

The energy question is the pivot on which the whole proposition turns, and it is worth being exact about who has measured what. Nobody has published a measured energy per bit for a TSV-stacked NAND stack read over UCIe; there is no such stack outside a fab. What exists is three kinds of evidence. At the bottom, the measured datacentre case: a Seoul National University analysis of streaming mixture-of-experts weights from commodity SSDs found the flash path costing about 102\,pJ/bit against HBM's 4.2, and concluded that flash needs a roughly tenfold improvement---to about 10\,pJ/bit---to be energy-competitive~\cite{ssd-energy} [V]. Commodity NVMe cannot do this job, and the appliance of Section~\ref{sec:appliance-cost} does not use it. In the middle, the vendors' relative statements: SanDisk's fact sheet says HBF's power profile is ``similar'' to HBM4's and that Gen.\ 2 and 3 draw 0.8 and 0.64 times Gen.\ 1; SK hynix's own H$^3$ paper, less generously, models HBF at up to four times HBM power at the same bandwidth~\cite{sandisk-hbf-facts,hbf-h3} [V]. At the top, the third-party models that assume the on-package figure: Ma and Patterson's ballpark of a 512\,GB stack at 1.6\,TB/s inside 80\,W, which is about 6\,pJ/bit; the Chinese Academy of Sciences' FlashAccel at 8\,pJ/bit; Micron's HotInfra roofline at 8.5~\cite{hbf-ma-patterson,hbf-flashaccel,hbf-micron-hotinfra} [V]. Those three converge on two to three times HBM per bit read, which is precisely the improvement the SNU paper says is required, and precisely the assumption Section~\ref{sec:appliance-cost} carries---7\,pJ/bit, 360\,W of flash reads at 6.4\,TB/s. The reader should hold that number as a modelled consensus, not a measurement, and should note that all three modelling groups are, in different degrees, parties to the outcome: Google is a consortium member, the CAS group is designing for it, and Micron is a competitor arguing that LPDDR wins.

Where the physics is settled is in what flash cannot do. A NAND page read takes tens of microseconds for triple-level cells and a few microseconds for single-level; the microsecond latency is intrinsic, and every HBF architecture study hides it rather than removes it---SK hynix's H$^3$ needs a 40\,MB SRAM latency-hiding buffer with a hit rate above 80\%, FlashAccel reads in $\sim$19\,MB ``hyper-pages'' behind SRAM prefetch, and Micron's model finds only 36\% of nominal bandwidth survives scattered access~\cite{hbf-h3,hbf-flashaccel,hbf-micron-hotinfra} [V]. Weights, read in a deterministic order decided by the model graph, are exactly the access pattern that tolerates this; John Carmack's remark that inference ``can have a deterministic memory access pattern \ldots could tolerate cold-start latencies in multiple milliseconds'' is the practitioner's version of the point~\cite{hbf-analysis} [A]. Key--value cache, written every token, is exactly the pattern that does not, and the August 2026 characterization by a Peking University group---HBF used for KV cache raising latency 2--5.5 times and cutting throughput 1.1--2.7 times---is the negative result everyone in the field expected~\cite{hbf-pku} [A; paper reported, not fetched]. And retention is a thermal problem before it is an endurance one: the HeatWatch study of 3D charge-trap NAND found retention loss accelerating superlinearly with temperature, which matters for a die stacked next to a kilowatt-class GPU and is one reason no vendor has yet published a temperature specification~\cite{hbf-heatwatch} [V]. The duty-cycle argument of Section~\ref{sec:memoryunlock}---that a non-volatile tier's zero standby draw offsets its higher active draw for a machine that is not busy around the clock---does not depend on any of the unmeasured numbers, and it is the part of the energy case this book is most confident of.

\section{Endurance: the number nobody has printed}

No vendor has published a program--erase cycle count, a terabytes-written figure, a read-disturb rate or a retention specification for HBF, and this appendix has confirmed the absence across every primary and secondary source it could reach. What exists is inference from the cell type and from the papers. If the cell is single-level, as multiple sources report, endurance on the order of 100,000 cycles is conventional and is what SK hynix's own H$^3$ paper assumes; FlashAccel assumes a million by relaxing retention, which is a design choice rather than a measurement; a triple-level cell would sit at three to eleven thousand~\cite{hbf-h3,hbf-flashaccel,hbf-semieng} [V/A]. For a weights tier that is written once per model load and read continuously, none of those figures is binding: an appliance that reloads a 1.4\,TB checkpoint even daily performs a few hundred full writes a year. Read disturb, the slow corruption of neighbouring cells by repeated reads, is the mechanism that actually applies to a read-mostly tier, and the reclaim writes it forces are wear; the STRAW work at SNU shows the disturbance is strongly non-uniform across word lines and that stress-aware reclaim cuts those writes by more than 80\%, which is the sort of controller-level engineering the base die exists to do~\cite{hbf-straw} [V]. SanDisk's own vice-president acknowledged the write-wear problem at FMS 2026 and said the risk ``diminishes if HBF operates primarily in read-heavy roles''~\cite{hbf-businesskorea} [A]. That is the correct statement of the case, and it is why the main text calls HBF a weights tier and never a memory replacement.

\section{Cost structure: cheap bits, expensive stack}

The cost argument is where the vendors' adjectives and the book's numbers must be reconciled, and the reconciliation is uncomfortable in both directions. Start with the bits. TrendForce's spot series for 512\,Gb TLC NAND wafers ran \$2.80 in August 2025, \$9.61 in December 2025 and \$21.13 on 9~August 2026---roughly \$0.044, \$0.15 and \$0.33 per gigabyte---so the raw die content of a 512\,GB stack was about \$22 at the pre-surge floor and about \$170 at today's spot~\cite{nand-spot} [V]. Custom 256\,Gb dies with through-silicon vias and possibly single-level cells cost several times commodity per bit, but the order of magnitude is unchanged: the die is not where the money is. The money is in the stack. TSV formation, thinning and stacking of sixteen or seventeen dies, the CBA base die with its sixteen host channels and their controllers, known-good-die test and burn-in are what SanDisk means when it says the stack costs about what an HBM stack costs, and what the Chipstrat analysis means when it puts HBF at roughly a tenth of HBM's cost per gigabyte but comparable per stack~\cite{sandisk-hbf-facts,hbf-analysis} [A]. An HBM4 12-high stack was reported at about \$500--560 in late 2025---\$14--16 per gigabyte---so ``similar cost'' would put a first-generation HBF stack near \$500, or about a dollar a gigabyte~\cite{hbf-hbm4-price} [A].

Section~\ref{sec:appliance-cost} prices the stack at \$150--250. The gap needs to be stated plainly rather than hidden. Three arguments carry it, in descending order of strength. First, HBM's \$500 is a \emph{price} in a market where HBM is the scarcest component in the industry and its vendors book record margins; it is not the cost of stacking sixteen dies. A finished 1\,TB NVMe SSD---controller, DRAM, PCB, enclosure, retail margin---sold for \$65--100 before the 2026 surge, and the operations that turn commodity NAND into an HBF stack are not obviously more expensive than the ones that turn it into a retail SSD, only different~\cite{nvme-presurge}. Second, the book's basis is declared to be pre-surge, and the vendors' ``similar to HBM'' was said in a market where everything memory-related has doubled or tripled since. Third, the appliance is a volume consumer product, and its stack does not compete with HBM for the same CoWoS interposer slots that a datacentre package's HBF would. Against those, the honest counter is that no HBF stack has been priced by anyone, that the vendors' one statement points higher, and that a two-vendor part in its first generation carries the margin a two-vendor part carries. So the sensitivity is given here rather than argued away: at \$500 per stack the appliance's base case rises from \$5,000 to about \$6,750 and its cost at 100 million tokens a month from \$1.65 to \$2.12 per million; at \$600, to \$7,300 and \$2.28 [M]. Neither moves the machine out of the order-of-magnitude gap against the closed API price list, and both move it above the hyperscaler's rack at planned utilization while leaving it below the rack at disappointed utilization. The stack price decides how much of the argument survives, not whether it does.

The market's own view of the cost structure is that it is unproven. Counterpoint's read of the August 2026 investor day was that SanDisk's own FY2028--30 financial model contains no HBF revenue---HBF is ``upside''; Wedbush declined to build it into estimates ``until adoption becomes clearer''; Mizuho's view is that it is ``unlikely to fully replace HBM''~\cite{hbf-counterpoint,hbf-tapeout} [A]. Market-size figures in circulation---\$1B in 2027 rising to \$12B in 2030 (unattributed, Korean press); ``bigger than HBM by 2038'' (a KAIST professor)---are not analyst forecasts and are cited here as what is being said, not as estimates~\cite{hbf-koreaherald,hbf-kaist} [A]. SanDisk's shares, up more than twenty-fold in a year on the NAND supercycle, rose 13.7\% on the investor day; the sell-side attributed most of that to the financial model and long-term agreements and only part to HBF~\cite{hbf-tapeout,hbf-counterpoint} [A].

\section{The perceived advantages, and the case against each}

The proponents' case is compact and this appendix states it at full strength before answering it. Capacity: 512\,GB per stack against 36--64\,GB for HBM4, so a whole 400-billion-to-3-trillion-parameter model sits on one package---SanDisk's ``4\,TB of VRAM,'' The Register's 3.12\,TB from two HBM and six HBF stacks, the H$^3$ paper's 32-GPU workload on two~\cite{hbf-toms-2025,hbf-register,hbf-h3}. Bandwidth: 1.6\,TB/s per stack matches HBM3E and the 3\,TB/s top grade exceeds a single HBM4 stack today, and decode is bandwidth-bound and read-mostly, which is the workload~\cite{hbf-spec,hbf-toms-spec}. Cost: an order of magnitude below HBM per gigabyte~\cite{hbf-analysis}. Non-volatility: no refresh, instant residency, model swap without reload from storage. Fewer accelerators for the same tokens: SanDisk's ``four HBF GPUs deliver the token output of eight HBM GPUs,'' Goldman's ``half the GPUs,'' FlashAccel's 2.54$\times$ throughput per GPU, H$^3$'s 2.69$\times$ per watt~\cite{hbf-investorday,hbf-tapeout,hbf-flashaccel,hbf-h3}. Two independent authorities give the concept cover: Patterson and Ma's research-directions paper proposes HBF as ``10$\times$ weight memory'' for mixtures of experts, and Carmack's determinism argument answers the latency objection~\cite{hbf-ma-patterson,hbf-analysis}. And the standardization moved fast: memorandum to consortium to specification in twelve months, with Google inside~\cite{hbf-ocp}.

The case against, item by item. \emph{The bandwidth comparison is against the wrong HBM.} SanDisk's investor-day slide set eight HBF stacks at 12.8\,TB/s against an HBM3E-era part; HBM4E stacks are specified near 3.6--4.1\,TB/s each, so an eight-stack HBM4E part in 2028 delivers roughly 30\,TB/s---HBF's read bandwidth is HBM's of two generations ago, and the gap widens on the roadmap~\cite{hbf-xeno,hbf-spectrum} [A]. \emph{Quantization shrinks the problem HBF solves.} The same slide sized Qwen3-480B at 960\,GB in BF16; at FP8 it is 480\,GB and at FP4---the format Kimi~K3 ships in natively---240\,GB, inside a large HBM4E complement; the appliance case of Section~\ref{sec:appliance-cost} is immune to this objection only because it targets a 2.8-trillion-parameter model at four bits, which is 1.4\,TB, and because it has no HBM at all~\cite{hbf-xeno} [A]. \emph{Latency and granularity are hidden, not solved}, and every study that hides them does so with SRAM and prefetch that a real base die must pay for~\cite{hbf-h3,hbf-flashaccel}. \emph{Key--value cache stays in DRAM}, so an HBF-only accelerator does not exist; the tier is additive to HBM or LPDDR, not a substitute, and Micron's roofline finds that once one is adding a tier anyway, on-package LPDDR---writable, 2.4\,pJ/bit---beats HBF at every operating point for datacentre decode by up to 2.1$\times$~\cite{hbf-micron-hotinfra} [V; a competitor's analysis]. \emph{Packaging is the cost and the beachfront.} If the stack costs what an HBM stack costs, the cheap bits buy capacity, not a cheap part; and an HBF stack on a datacentre package occupies an interposer slot an HBM stack would otherwise take~\cite{hbf-analysis}. \emph{Two vendors, and the wrong absentees.} NVIDIA has said nothing; AMD, Micron and Samsung are outside; Kioxia is building something else with NVIDIA; the trade press's summary in January 2026 was ``get Nvidia on board \ldots and it's a done deal,'' and the corollary holds~\cite{hbf-bf-window,hbf-toms-spec} [A]. \emph{Optane.} A proprietary memory tier at two vendors, needing a new controller, a new interface and a new software layer, has failed before on adoption rather than physics~\cite{hbf-optane-analogy} [A]. \emph{Timelines.} Samples have moved from 2H26 to 2027 in twelve months, production to 2028, and SK hynix has given no production date at all~\cite{hbf-techradar,hbf-tapeout} [A]. \emph{Evidence.} Every performance figure in circulation is a simulation; the die taped out in August 2026 has not been characterized by anyone outside SanDisk~\cite{hbf-tradethepool} [A].

The book's reading of that exchange is the one Section~\ref{sec:memoryunlock} gives: the case against is strong precisely where the proponents overclaim---HBF as HBM's rival in the datacentre package---and weak where the book actually uses it---HBF as the weights tier of a machine that has no HBM, serves a handful of users, holds a four-bit model larger than any HBM complement, and is idle two-thirds of the day. Micron's roofline, the sharpest of the criticisms, is a datacentre-decode result at high batch with a writable tier competing for the same slot; the appliance's KV tier is LPDDR6X, exactly as Micron would prefer, and its weights tier is chosen for capacity per dollar rather than bandwidth per watt. Whether the appliance's flash reads cost 6 or 8.5 or 12\,pJ/bit changes its power draw by a hundred watts and its running cost by a few dollars a month. What would change the argument is HBF failing to ship, or shipping at a stack price closer to \$1,000 than \$200---and both are live possibilities that the forecast register carries.

\section{The competitive landscape}

HBF's rivals for the ``tier between HBM and storage'' are of three kinds, and the book's position on each is stated so it can be checked. \emph{More HBM}: HBM4 at 36--48\,GB per stack with a JEDEC ceiling of 64\,GB, HBM4E sampled in May 2026 for late-2027 volume at 3.6--4.1\,TB/s---which addresses bandwidth handsomely and capacity hardly at all, and at the highest cost per bit in the industry~\cite{hbf-hbm-roadmap,hbf-samsung-fms26} [A]. \emph{Capacity DRAM}: LPDDR5X SOCAMM2 modules at 256\,GB (2\,TB per server) already shipping, LPDDR6 on JEDEC's data-centre roadmap, Qualcomm's AI200 with 768\,GB of LPDDR per card, CXL memory expansion, and Micron's argument that this tier dominates HBF for datacentre decode; the appliance uses it for context and would use it for weights too if flash did not arrive~\cite{hbf-micron-hotinfra,jedec-lpddr6-roadmap,qualcomm-ai200} [V/A]. \emph{Other flash}: Kioxia's path---the 5\,TB, 64\,GB/s daisy-chained module of August 2025, a 100-million-IOPS GPU-direct SSD for 2027 built with NVIDIA, XL-Flash CXL modules, and a concept of up to 32 TSV-stacked dies on a shared interposer beside HBM~\cite{hbf-kioxia-module,hbf-kioxia-tsv} [V/A]; Samsung's zNAND-O, a four-to-eight-die NAND stack on a TSV base die with UCIe to an on-device NPU, targeted at 2028---HBF's edge concept, executed outside the OCP tent~\cite{hbf-samsung-fms26} [V]; and the research thread that puts compute \emph{inside} the flash rather than beside it (AiF at ISCA 2025, KVNAND, NVLLM, and a Korea University 3D-NAND processing-in-memory array), which is where a Chinese ecosystem with a strong NAND base and weak HBM could plausibly go~\cite{hbf-inflash} [V]. On the DRAM side, ZAM (Section~\ref{sec:memoryunlock}) is the capacity-by-height alternative that would compete with HBF for the same weights-tier role at DRAM latency, later and at a higher price per bit.

The China angle is worth stating separately because it is not what the Western coverage assumes. Chinese analysts have framed HBF as a ``lane-change'' opportunity: China ``does not have high-end production capability in HBM,'' but YMTC's Xtacking is a wafer-bonded architecture of exactly the CBA kind HBF needs, YMTC's 294-layer part is reported in production at yields above 90\% with a 13\% unit share and an expansion from 200 to 300 thousand wafers a month by 2027, and its hybrid-bonding intellectual property has been licensed to Samsung~\cite{hbf-huxiu} [A]. YMTC's own developer conference in December 2025 said HBF ``fits Xtacking'' and DigiTimes reported it going ``full speed toward HBF next''; a competing report has YMTC prioritizing HBM with CXMT instead~\cite{hbf-ymtc,hbf-crazyrich} [A]. The constraint is tooling---the entity list bars the equipment for beyond roughly 300 layers---and the decider, in the Chinese framing as in the Western one, is whether NVIDIA adopts the tier. The book's Chapter~\ref{ch:semi} argument that a NAND-rich, HBM-poor ecosystem is well placed for a flash weights tier is consistent with that framing but is not yet supported by a Chinese HBF product announcement, and it is graded accordingly.

\section{The research literature}

Table~\ref{tab:hbf-papers} lists the papers this appendix relies on and the ones a reader should know exist. Three strands run through it. The \emph{offloading} strand, from FlexGen and Apple's ``LLM in a flash'' in 2023 through the mixture-of-experts expert-caching work of 2024--26 (MoE-Infinity, MoE-Lightning, ExpertFlow, PowerInfer-2 on a phone), established that weights can be streamed from flash behind a small hot cache and that expert routing has enough locality to make caching worthwhile, but that the SSD path is too slow and, per the SNU paper, too energy-hungry to be more than a fallback~\cite{hbf-offload,ssd-energy}. The \emph{HBF-specific} strand, all from 2026, models the on-package tier and agrees on its shape: SK hynix's H$^3$, the POSTECH--ETH characterization (which concludes HBF helps ``only if HBM-equivalent read performance and substantial endurance improvement are achieved''), FlashAccel, Ma and Patterson, and Micron's dissent~\cite{hbf-h3,hbf-postech,hbf-flashaccel,hbf-ma-patterson,hbf-micron-hotinfra}. The \emph{routing-locality} strand is what decides whether a flash weights tier serves cost or bandwidth: per-token sparsity is high---DeepSeek-R1 activates 3.5\% of parameters---but batching erases it, so that a batch of a few dozen touches nearly every expert; local routing consistency varies widely across models and is degraded by shared-expert designs; a hot cache of about twice the active experts captures most of the locality that exists~\cite{hbf-routing,ssd-energy}. That last strand is the literature's version of the batching-saturation arithmetic of Section~\ref{sec:memoryunlock}, and it cuts the same way: an ultra-sparse mixture of experts is exactly the model for which a capacity tier is worth having and a bandwidth-amortizing rack is worth less.

\begin{table}[htbp]
  \centering
  \scriptsize
  \caption{Selected literature on flash weights tiers and HBF. All arXiv identifiers verified; venue given where published.}
  \label{tab:hbf-papers}
  \begin{tabular}{p{42mm}p{26mm}p{70mm}}
    \toprule
    Paper & Venue / date & Result relevant to this book \\
    \midrule
    FlexGen (Sheng et al.) & ICML 2023 & OPT-175B on one 16\,GB GPU with disk offload at $\sim$1 token/s: the floor \\
    LLM in a flash (Apple) & ACL 2024 & windowing + row--column bundling: models 2$\times$ DRAM, 20--25$\times$ over naive flash loading \\
    Fast MoE inference with offloading (Eliseev, Mazur) & Dec.\ 2023 & Mixtral: LRU expert cache hit 0.4--0.8; 2--3 tok/s on consumer GPUs \\
    MoE-Infinity; MoE-Lightning; ExpertFlow; HOBBIT; PreScope; FlashMoE & 2024--26 & expert caching, prefetch and prediction: 3--16$\times$ latency and up to 10$\times$ throughput over naive offload; SSD-resident experts on the edge \\
    PowerInfer-2 (SJTU) & Jun.\ 2024 & 47B MoE streamed from smartphone UFS flash at 11.7 tok/s \\
    SSD offloading considered harmful (SNU) & CAL 2025 & flash path $\sim$102\,pJ/bit vs HBM 4.2; per-token energy 3.8--12.5$\times$; needs $\sim$10\,pJ/bit \\
    Ma \& Patterson & arXiv 2601.05047, Jan.\ 2026 & HBF as 10$\times$ weight memory; stack ballpark 512\,GB, 1.6\,TB/s, $<$80\,W, 4\,KB granule \\
    H$^3$ (SK hynix) & IEEE CAL, Jan.\ 2026 & 8 HBM + 8 HBF: 2.69$\times$ perf/W, 6.14$\times$ throughput at 10M context; 40\,MB latency buffer; $\sim$100k cycles \\
    Exploring HBF for LLM inference (POSTECH, ETH) & IEEE CAL, Jan.\ 2026 & gains ``only if HBM-equivalent read performance and substantial endurance improvement'' \\
    HAVEN (UCSD, Georgia Tech) & arXiv 2603.01175, Mar.\ 2026 & HBF for full-precision vector search in RAG: up to 20$\times$ throughput \\
    Micron, Is HBF all you need? & HotInfra at ISCA, Jun.\ 2026 & on-package LPDDR beats HBF at every point for datacentre decode; HBF 8.5\,pJ/bit, 36\% scattered-read efficiency \\
    FlashAccel (ICT-CAS) & arXiv 2607.10186, Jul.\ 2026 & 6 HBF stacks + HBM: 2.54$\times$ throughput/GPU, 1.93$\times$ tokens/J; 8\,pJ/bit, 4\,\textmu s tR, 19\,MB hyper-pages \\
    HBF for KV-centric serving (Peking Univ.) & arXiv, Aug.\ 2026 (reported) & KV cache on HBF: latency 2--5.5$\times$, throughput 1.1--2.7$\times$ worse \\
    AiF (SNU, ISCA 2025); KVNAND; NVLLM; 3D-NAND PIM (Korea Univ.) & 2025--26 & compute inside the flash: 14.6$\times$ over SSD offload; DRAM-free edge LLM at $\sim$65 tok/s on 30B \\
    Local routing consistency (Fudan, Huawei); Opportunistic expert activation & 2025--26 & routing locality varies across 20 MoE models; cache $\approx$2$\times$ active experts near-optimal; batch-aware rerouting cuts MoE decode latency 15--39\% \\
    HeatWatch (CMU, ETH); STRAW (SNU) & HPCA 2018; CAL 2024 & 3D NAND retention superlinear in temperature; read-disturb reclaim writes $-$84\% with word-line-aware policy \\
    \bottomrule
  \end{tabular}
\end{table}

\section{What the main text assumes, and what would falsify it}

Section~\ref{sec:appliance-cost} assumes four first-generation stacks at 512\,GB and 1.6\,TB/s each, 7\,pJ/bit read energy, zero standby power, a stack price of \$150--250 on a pre-surge basis, and endurance adequate for a device that reloads its model at most daily. Against the record in this appendix: the capacity and bandwidth are the vendors' published Gen.\ 1 figures and the specification's middle grade [V, as targets]; the energy is the third-party modelled consensus and about twice HBM's [M]; the standby figure is the fact sheet's non-volatility claim [V]; the price is below the vendors' one adjective and is the assumption the appliance's cost is most sensitive to [M]; the endurance is unpublished but benign for the workload [M]. The forecast register (Appendix~\ref{app:dashboard}) carries the resolving events: HBF-class silicon in volume at $\geq$512\,GB and $\geq$1.6\,TB/s with a published endurance figure; the appliance configuration below \$8k. To those this appendix adds the four observations that would move the argument fastest, in either direction: an independently measured energy per bit for a TSV-stacked NAND stack read over UCIe; a published endurance and retention specification at package temperature; a per-stack price, from anyone, against HBM4E's; and a design-in by NVIDIA, AMD or a hyperscale ASIC---or, equally informative, Kioxia's and Samsung's separate paths reaching product first. As of \datafreeze\ none of the four has happened, one die has taped out, and the argument of Chapter~\ref{ch:compute} stands on a specification, a fact sheet, three simulations and a duty-cycle calculation. That is more than it stood on a year ago and less than it will need.
